\chapter{模块详细设计}

\section{身份认证模块}
身份认证模块是整个 LingoBridge 系统的安全基石，负责管理用户从匿名游客到注册用户的状态跃迁。前端的核心逻辑封装在全局状态管理器中，提供用于分发用户状态以及身份令牌的组件。对于未登录的用户，系统将统一挂载一个标识，使得基础导航和公开内容对游客开放。而对于需要高级操作权限的按钮、输入框或特定页面，系统设计了高级守卫组件。该组件接收子节点并监控当前用户的角色，一旦检测到用户试图触发未授权操作，它将自动拦截点击事件并拉起居中浮动的动态登录弹框。当用户输入凭证并调用鉴权接口取得令牌后，信息会被同时存储于客户端的本地持久化缓存与服务请求头中。后端的请求适配器会自动提取该令牌并进行解密与时效校验，若校验失败则立即向上层抛出未授权状态，触发客户端清除缓存并重定向至首页。

\section{课程管理模块}
课程管理模块由教师端核心视图组件构成，提供多层次的课程和课时维护能力。在教师的课程列表主页中，系统采用响应式卡片网格渲染教师所创建的全部课程，并支持关键字搜索与类型过滤。教师点击任意课程卡片后，将无缝载入课时详情视图组件。该组件采用树形嵌套结构展示各个教学单元、具体课时节点以及每个课时下挂载的作业评估任务。为了支持高频的创建、更新与删除操作，前端引入了数据流状态，将网络请求返回的原始数组转化为具有高度响应性的内存树状结构。当教师点击添加课时时，组件将拉起滑出式侧边面板并收集课时标题、预计时长与难度等级，随后通过底层的服务代理发送结构化请求。后端接收后会同步生成具有全局唯一标识符的课时节点，同时为之预初始化一个一对一绑定的作业节点，以简化后续学生的作业提交路径。

\section{课件渲染模块}
课件渲染模块需要解决教师上传的演示文稿与便携式文档在不同设备屏幕上的高保真还原问题。系统采用基于网格定位的层叠渲染架构。底层是利用第三方开源渲染引擎将上传至服务器的文档渲染到高清晰度的物理画布上。为了防止在高分辨率视网膜屏幕上出现图像模糊现象，渲染器会动态读取设备的物理像素比，并在初始化画布时对其物理尺寸进行等比例放大，最后通过样式宽度控制其显示区域，从而实现高保真度显示。在文档画布之上的正上方，系统通过绝对定位叠加了一个完全透明的交互式绘图层。该绘图层拦截一切鼠标和触控手势输入，使得用户的滑动手势能够实时转化为屏幕坐标系中的平滑贝塞尔曲线。为了确保绘制的线条在各种分辨率与拉伸比例下均不变形，系统将所有绘制动作转换为归一化坐标保存，实现了设备无关的涂鸦展示。

\section{录音与测评模块}
录音与语音评测模块旨在为学生提供即时的发音反馈。在前端，模块深度整合了浏览器的媒体录制接口，在组件挂载时自动申请用户的麦克风使用授权，并构建了防抖机制以避免因频繁点击造成的音频流锁死。一旦录音开始，系统将每隔一定毫秒切片捕获音频数据，将其作为二进制缓冲区存入本地数组中。录音结束后，组件会自动将这组缓冲区合并编码，包装成标准的无损波形音频文件格式，并通过专门的文件传输通道向后端提交。为了提升响应速度并支持长录音的高效传输，系统在高级网络协议上建立了音频传输套接字，支持将分片缓冲区以二进制流的形式低延迟泵送至后端。后端在接收到完整的音频文件后，会将其安全写入文件存储区，并在数据库记录中关联该音频的路径，随后将音频特征发送至发音评估引擎，并将评分与音素级别的错误纠错信息异步返回给学生终端。

\section{直播课堂模块}
直播课堂模块是促进师生双向实时互动的高频应用场景。该模块包含了课件同步、实时电子白板涂鸦、音视频通话以及文字聊天四个子模块。为了保证教师端与多名学生端之间的绝对一致性，系统采用主备控制模型。教师端作为整个课堂的决策主导者，其当前查看的课件页码、电子白板的笔迹更新以及清理画布等操作均会被提炼为高层语义事件，通过持久连接层广播给所有处于活跃状态的学生终端。电子白板涂鸦的状态管理以页码为唯一索引，前端在内存中维护着多页的笔迹状态树。当教师或学生在某一页上绘制线条时，渲染器仅更新该页的本地画布上下文。若触发清理画布操作，系统会识别并仅清空当前活动页对应的笔迹数组，而不会误伤其他未显示页面的历史记录。对于学生的白板绘制权限，教师端可通过权限控制面板随时开启或禁用，确保直播教学的有序进行。
